\documentclass{elsarticle}
\usepackage[utf8]{inputenc}
\usepackage[T1]{fontenc}
\usepackage{geometry}
\usepackage{graphicx}
\usepackage{float}
\usepackage{hyperref}
\usepackage[english,indonesian]{babel}

\title{Optimasi Pemrosesan Template LaTeX Menggunakan Pendekatan Berbasis AI}
\author{Academic Support Agent}
\date{\today}

\begin{document}
\maketitle

\begin{abstract}

\end{abstract}

\chapter{PENDAHULUAN}


\chapter{TINJAUAN PUSTAKA}


% Original content (commented out):
% % Template for Elsevier CRC journal article
% version 1.1-5p dated 18 January 2011
% SEM ACENTUACAO PROBLEMA UTF-8 : NAO ACEITA

% This file (c) 2010-2011 Elsevier Ltd.  Modifications may be freely made,
% provided the edited file is saved under a different name

% This file contains modifications for Procedia Computer Science
% but may easily be adapted to other journals

% Changes since version 1.0
% - elsarticle class option changed from 1p to 3p (to better reflect CRC layout)
% - this version uses option 5p for larger-format journals (text area 24.1 x 18.4 cm)
% Newest modifications by Antef Nava. Dated 5 November 2016
% Now accepts accents and can add only two logos, this modifications is in the sty file
%-----------------------------------------------------------------------------------

%% This template uses the elsarticle.cls document class and the extension package ecrc.sty
%% For full documentation on usage of elsarticle.cls, consult the documentation "elsdoc.pdf"
%% Further resources available at http://www.elsevier.com/latex

%-----------------------------------------------------------------------------------

%%%%%%%%%%%%%%%%%%%%%%%%%%%%%%%%%%%%%%%%%%%%%%
%%%%%%%%%%%%%%%%%%%%%%%%%%%%%%%%%%%%%%%%%%%%%%
%%                                          %%
%% Important note on usage                  %%
%% -----------------------                  %%
%% This file must be compiled with PDFLaTeX %%
%% Using standard LaTeX will not work!      %%
%%                                          %%
%%%%%%%%%%%%%%%%%%%%%%%%%%%%%%%%%%%%%%%%%%%%%%
%%%%%%%%%%%%%%%%%%%%%%%%%%%%%%%%%%%%%%%%%%%%%%

%% The '5p' and 'times' class options of elsarticle are used for Elsevier CRC
\documentclass[5p,times,authoryear]{elsarticle}

%% The `ecrc' package must be called to make the CRC functionality available
%% ecrc_RIAI es el paquete ecrc de Elsevier con modificaciones para la revista RIAI
\usepackage{ecrc_RIAI}

%% The ecrc package defines commands needed for running heads and logos.
%% For running heads, you can set the journal name, the volume, the starting page and the authors

%%%%%%%%%%%%%%%%%%%%%%%%%%%%%%%%% Aaadido por Secretaraa RIAI
\usepackage[spanish]{babel}     % Idioma
\usepackage[utf8x]{inputenc}
\usepackage[spanish]{layout}
\addto\captionsspanish{%
\def\tablename{Tabla}%
}

\usepackage{amsmath}            % Para las referencias a ecuaciones con \eqref
\usepackage{epstopdf}           % Para poder insertar figuras .eps al compilar con PDFLATEX
\usepackage{flushend}           % Para igualar las columnas de la altima pagina
%\usepackage{hyperref}           % Para hipervanculos dentro del PDF
%%%%%%%%%%%%%%%%%%%%%%%%%%%%%%%%%%%%%%%%%%%%%%%%%%%%%%%


%% set the volume if you know. Otherwise `00'
\volume{00}

%% set the starting page if not 1
\firstpage{1}

%% Give the name of the journal
\journalname{Prensas de Ciencias}

%% Give the author list to appear in the running head
%% Example \runauth{C.V. Radhakrishnan et al.}
\runauth{Primer autor et al.}

%% The choice of journal logo is determined by the \jid and \jnltitlelogo commands.
%% A user-supplied logo with the name <\jid>logo.pdf will be inserted if present.
%% e.g. if \jid{yspmi} the system will look for a file yspmilogo.pdf
%% Otherwise the content of \jnltitlelogo will be set between horizontal lines as a default logo

%% Give the abbreviation of the Journal. Contast the Publisher if in doubt what this is.
\jid{RIAI}

%% Give a short journal name for the dummy logo (if needed)
\jnltitlelogo{}

%% Hereafter the template follows `elsarticle'.
%% For more details see the existing template files elsarticle-template-harv.tex and elsarticle-template-num.tex.

%% Elsevier CRC generally uses a numbered reference style
%% For this, the conventions of elsarticle-template-num.tex should be followed (included below)
%% If using BibTeX, use the style file elsarticle-num.bst

%% End of ecrc-specific commands
%%%%%%%%%%%%%%%%%%%%%%%%%%%%%%%%%%%%%%%%%%%%%%%%%%%%%%%%%%%%%%%%%%%%%%%%%%

%% The amssymb package provides various useful mathematical symbols
\usepackage{amssymb}
%% The amsthm package provides extended theorem environments
%% \usepackage{amsthm}

%% The lineno packages adds line numbers. Start line numbering with
%% \begin{linenumbers}, end it with \end{linenumbers}. Or switch it on
%% for the whole article with \linenumbers after \end{frontmatter}.
%% \usepackage{lineno}

%% natbib.sty is loaded by default. However, natbib options can be
%% provided with \biboptions{...} command. Following options are
%% valid:

%%   round  -  round parentheses are used (default)
%%   square -  square brackets are used   [option]
%%   curly  -  curly braces are used      {option}
%%   angle  -  angle brackets are used    <option>
%%   semicolon  -  multiple citations separated by semi-colon
%%   colon  - same as semicolon, an earlier confusion
%%   comma  -  separated by comma
%%   numbers-  selects numerical citations
%%   super  -  numerical citations as superscripts
%%   sort   -  sorts multiple citations according to order in ref. list
%%   sort&compress   -  like sort, but also compresses numerical citations
%%   compress - compresses without sorting
%%
%% \biboptions{comma,round}

% \biboptions{}

% if you have landscape tables
\usepackage[figuresright]{rotating}

% put your own definitions here:
%   \newcommand{\cZ}{\cal{Z}}
%   \newtheorem{def}{Definition}[section]
%   ...

% add words to TeX's hyphenation exception list
%\hyphenation{author another created financial paper re-commend-ed Post-Script}

% para poder introducir varias figuras que ocupen el ancho de las dos columnas.
\usepackage{subfigure}

% declarations for front matter

\begin{document}

\begin{frontmatter}

%% Title, authors and addresses

%% use the tnoteref command within \title for footnotes;
%% use the tnotetext command for the associated footnote;

%% use the fnref command within \author or \address for footnotes;
%% use the fntext command for the associated footnote;

%% use the corref command within \author for corresponding author footnotes;
%% use the cortext command for the associated footnote;
%% use the ead command for the email address,
%% and the form \ead[url] for the home page:
%%
%% \title{Title\tnoteref{label1}}
%% \tnotetext[label1]{}
%% \author{Name\corref{cor1}\fnref{label2}}
%% \ead{email address}
%% \ead[url]{home page}
%% \fntext[label2]{}
%% \cortext[cor1]{}
%% \address{Address\fnref{label3}}
%% \fntext[label3]{}

%\dochead{Cabecera artaculo}
%% Use \dochead if there is an article header, e.g. \dochead{Short communication}
%% use optional labels to link authors explicitly to addresses:
%% \author[label1,label2]{<author name>}
%% \address[label1]{<address>}
%% \address[label2]{<address>}

\title{Plantilla simplificada para la preparación \\
de artículos en Overleaf}

\author[First]{Primer A. Autor\corref{cor1}} %\fnref{label2}} %Nota al pie
\ead{autor@cea-ifac.es}

\author[Second]{Segundo B. Autor, Jr.}
\ead{autor2@cea-ifac.es}

\author[Third]{Tercer C. Autor}
\ead{autor3@cea-ifac.es}

%\fntext[label2]{Nota al pie para el autor 1}
\cortext[cor1]{Autor en correspondencia.}

\address[First,Second,Third]{Facultad de Ciencias, Universidad Nacional Autónoma de México, Ciudad Universitaria, c.p. 04510, Ciudad de México}

%se pueden agregar más direcciones, descomenta la siguiente línea
%\address[Second]{Facultad de Ciencias, Universidad Nacional Autónoma de México, Ciudad Universitaria, c.p. 04510, Ciudad de México}

%resumn
\begin{abstract}
En el siguiente documento se puede observar como agregar secciones, imágenes, tablas, bibliografía y apéndices. \emph{Plantilla original por Alvarez, con licencia de uso Creative Commons CC BY 4.0. Adaptación para el uso de la clase de Introducción a la Geodinámica, Facultad de Ciencias, UNAM}
\end{abstract}

\begin{keyword}

palabra 1 \sep palabra 2 \sep 5-10 palabras clave (tomadas de la lista del sitio web de IFAC).

\end{keyword}

\end{frontmatter}

\tableofcontents

%LOS LETREROS DE SECCIÓN Y SUBSECCIÓN SÓLO SON UNA AYUDA VISUAL

%%============================= S E C C I Ó N =============================%%
\section{Introducción}
Las secciones, subsecciones e imágenes de este documento provienen del artículo de \cite{Ferrari2}

%%============================= S E C C I Ó N =============================%%
\section{Marco geológico}
	Lorem ipsum dolor sit amet, consectetur adipiscing elit. Nullam aliquet eros eu augue lobortis varius.
    
%%--------------------------- s u b s e c c i ó n ---------------------------%%
	\subsection{Introducción}
    Lorem ipsum dolor sit amet, consectetur adipiscing elit. Nullam aliquet eros eu augue lobortis varius.
    
%%--------------------------- s u b s e c c i ó n ---------------------------%%
	\subsection{Espesor y naturaleza de la placa superior}
        Lorem ipsum dolor sit amet, consectetur adipiscing elit. Nullam aliquet eros eu augue lobortis varius.
        
%%--------------------------- s u b s e c c i ó n ---------------------------%%
	\subsection{Geometría de la placa en subducción}
        Lorem ipsum dolor sit amet, consectetur adipiscing elit. Nullam aliquet eros eu augue lobortis varius.
        
%%--------------------------- s u b s e c c i ó n ---------------------------%%
	\subsection{Fluidos y fundidos en la cuña del manto y la corteza}
        Lorem ipsum dolor sit amet, consectetur adipiscing elit. Nullam aliquet eros eu augue lobortis varius.
        
%%============================= S E C C I Ó N =============================%%
\section{Evolución geológica y geoquímica}
    Lorem ipsum dolor sit amet, consectetur adipiscing elit. Nullam aliquet eros eu augue lobortis varius.
    
%%--------------------------- s u b s e c c i ó n ---------------------------%%
	\subsection{Distribución espacio temporal del vulcanismo}
        Lorem ipsum dolor sit amet, consectetur adipiscing elit. Nullam aliquet eros eu augue lobortis varius.
        
%%============================= S E C C I Ó N =============================%%
\section{Tectónica en el Neógeno}
    Lorem ipsum dolor sit amet, consectetur adipiscing elit. Nullam aliquet eros eu augue lobortis varius.
    
%%--------------------------- s u b s e c c i ó n ---------------------------%%
	\subsection{Introducción}
        Lorem ipsum dolor sit amet, consectetur adipiscing elit. Nullam aliquet eros eu augue lobortis varius.
        
%%--------------------------- s u b s e c c i ó n ---------------------------%%
    \subsection{Este de Jalisco, Michoacán y Guanajuato}
        Lorem ipsum dolor sit amet, consectetur adipiscing elit. Nullam aliquet eros eu augue lobortis varius.
        
%%--------------------------- s u b s e c c i ó n ---------------------------%%
    \subsection{Migración de la actividad sísmica}
        Lorem ipsum dolor sit amet, consectetur adipiscing elit. Nullam aliquet eros eu augue lobortis varius.
        
%%============================= S E C C I Ó N =============================%%
\section{Discusión}
    Lorem ipsum dolor sit amet, consectetur adipiscing elit. Nullam aliquet eros eu augue lobortis varius.
    
%%--------------------------- s u b s e c c i ó n ---------------------------%%
	\subsection{Causa de la subducción plana}
        Lorem ipsum dolor sit amet, consectetur adipiscing elit. Nullam aliquet eros eu augue lobortis varius.
        
%%--------------------------- s u b s e c c i ó n ---------------------------%%
    \subsection{Causa de la deformación intra arco}
        Lorem ipsum dolor sit amet, consectetur adipiscing elit. Nullam aliquet eros eu augue lobortis varius.
        
%%--------------------------- s u b s e c c i ó n ---------------------------%%
    \subsection{Observaciones importantes y preguntas pendientes}
        Lorem ipsum dolor sit amet, consectetur adipiscing elit. Nullam aliquet eros eu augue lobortis varius.

\begin{figure}
\centering
  \includegraphics[width=7cm]{ferrari}\\
  \caption{Aquí tienen que escribir el pie de figura}
\end{figure}





\bibliographystyle{elsarticle-harv}
\bibliography{riaibib}


\end{document}

%%
%% End of file 

\end{document}